% Options for packages loaded elsewhere
\PassOptionsToPackage{unicode}{hyperref}
\PassOptionsToPackage{hyphens}{url}
\PassOptionsToPackage{dvipsnames,svgnames,x11names}{xcolor}
\documentclass[
]{article}
\usepackage{xcolor}
\usepackage[margin=25mm]{geometry}
\usepackage{amsmath,amssymb}
\setcounter{secnumdepth}{-\maxdimen} % remove section numbering
\usepackage{iftex}
\ifPDFTeX
  \usepackage[T1]{fontenc}
  \usepackage[utf8]{inputenc}
  \usepackage{textcomp} % provide euro and other symbols
\else % if luatex or xetex
  \usepackage{unicode-math} % this also loads fontspec
  \defaultfontfeatures{Scale=MatchLowercase}
  \defaultfontfeatures[\rmfamily]{Ligatures=TeX,Scale=1}
\fi
\usepackage{lmodern}
\ifPDFTeX\else
  % xetex/luatex font selection
\fi
% Use upquote if available, for straight quotes in verbatim environments
\IfFileExists{upquote.sty}{\usepackage{upquote}}{}
\IfFileExists{microtype.sty}{% use microtype if available
  \usepackage[]{microtype}
  \UseMicrotypeSet[protrusion]{basicmath} % disable protrusion for tt fonts
}{}
\makeatletter
\@ifundefined{KOMAClassName}{% if non-KOMA class
  \IfFileExists{parskip.sty}{%
    \usepackage{parskip}
  }{% else
    \setlength{\parindent}{0pt}
    \setlength{\parskip}{6pt plus 2pt minus 1pt}}
}{% if KOMA class
  \KOMAoptions{parskip=half}}
\makeatother
\setlength{\emergencystretch}{3em} % prevent overfull lines
\providecommand{\tightlist}{%
  \setlength{\itemsep}{0pt}\setlength{\parskip}{0pt}}
\usepackage{bookmark}
\IfFileExists{xurl.sty}{\usepackage{xurl}}{} % add URL line breaks if available
\urlstyle{same}
\hypersetup{
  colorlinks=true,
  linkcolor={Maroon},
  filecolor={Maroon},
  citecolor={Blue},
  urlcolor={Blue},
  pdfcreator={LaTeX via pandoc}}

\author{}
\date{}

\begin{document}

\section{Generic zeros and absolute
primality}\label{generic-zeros-and-absolute-primality}

\emph{Written by GPT-6.1 Sol (OpenAI) in Codex at Ultra reasoning
effort, October 2026. Self-checked by the writing AI, GPT-6.1 Sol, at
Ultra. No independent review is claimed. Public domain (CC0).}

A prime ideal can describe a whole variety through one point, provided
that point is allowed to have transcendental coordinates. Its
specializations recover all the ordinary points. Changing the
coefficient field asks a further question: does this variety remain
integral, or does it split or acquire nilpotents? Generic zeros connect
these questions to field extensions.

Let \(k\) be any field, \(R=k[x_1,\ldots,x_n]\), and let
\(\Omega\supset k\) be algebraically closed. When generic points must
belong to \(\Omega^n\), assume \(\operatorname{trdeg}_k\Omega\ge n\).
For arbitrary points we may use different extension fields. We assume
the \href{../prerequisites/AG-CA/AG-CA-06.html}{Nullstellensatz},
Theorems 2.1--2.2, and
\href{../prerequisites/AG-CA/AG-CA-09.html}{dimension and Noether
normalization}, Theorems 4.2--4.3 and 5.1. The primary-decomposition
facts come from the preceding lessons and the associated-prime
prerequisite. The field-extension criteria in Section 5 are proved here.
Noether's working English edition and the AI Integrated Stacks Project
provide further reading.

\subsection{1. A point that remembers exactly an
ideal}\label{a-point-that-remembers-exactly-an-ideal}

For a tuple \(\xi\) in an extension field define

\[
I_k(\xi)=\{f\in R:f(\xi)=0\}.
\]

It is prime because evaluation has values in a field.

\textbf{Theorem 1.1 (generic zero).} For a prime
\(\mathfrak p\subset R\), put \(A=R/\mathfrak p\),
\(L=\operatorname{Frac}(A)\), and let \(\xi_i\) be the image of \(x_i\)
in \(L\). Then \(I_k(\xi)=\mathfrak p\).

\textbf{Proof.} Evaluation is the quotient map \(R\to A\) followed by
the injective map \(A\to L\). Its kernel is precisely \(\mathfrak p\).
\(\square\)

The tuple is called a \textbf{generic zero}. The field
\(L=k(\xi_1,\ldots,\xi_n)\) is its function field. It can be embedded
over \(k\) into the stated large \(\Omega\): send a transcendence basis
of \(L/k\), of size at most \(n\), to algebraically independent elements
of \(\Omega\), then extend over the finite algebraic extension to an
embedding in the algebraically closed field. The relation ideal stays
\(\mathfrak p\).

For \(\mathfrak p=(y-x^2)\), the generic zero is \((t,t^2)\), with \(t\)
transcendental. Evaluation gives \(k[x,y]/\mathfrak p\cong k[t]\), so no
additional relation is hidden.

\subsection{2. Specialization acts on coordinate
rings}\label{specialization-acts-on-coordinate-rings}

Say that \(\eta\) is a \textbf{specialization} of \(\xi\) if
\(I_k(\xi)\subset I_k(\eta)\). This definition preserves equations
rather than prescribing a limiting process.

\textbf{Theorem 2.1.} The condition is equivalent to a \(k\)-algebra
homomorphism

\[
k[\xi_1,\ldots,\xi_n]\longrightarrow k[\eta_1,\ldots,\eta_n],
\qquad \xi_i\longmapsto\eta_i.
\]

For a fixed target field \(F\supset k\), the specializations in \(F^n\)
are exactly \(V_F(I_k(\xi))\).

\textbf{Proof.} Both coordinate algebras are quotients of \(R\) by their
relation ideals. A map induced by the prescribed assignment exists
exactly when every relation of \(\xi\) is a relation of \(\eta\). The
last assertion restates this condition as evaluation of all polynomials
in the relation ideal. \(\square\)

Such a map is not in general a map of function fields. For example,
\(t\mapsto0\) defines \(k[t]\to k\) but cannot extend to \(k(t)\), since
\(t^{-1}\) would require an inverse of zero. This is why specialization
uses \(k[\xi]\).

The prime-inclusion picture of specialization in the spectrum is
precisely the same order: a larger prime imposes more equations. Closed
points have coordinates algebraic over \(k\), by the weak
Nullstellensatz; generic zeros may require transcendental coordinates.

\subsection{3. All zeros, including those of primary
ideals}\label{all-zeros-including-those-of-primary-ideals}

\textbf{Proposition 3.1.} For every ideal \(I\), \(V_F(I)=V_F(\sqrt I)\)
in every extension field \(F\). If \(Q\) is \(\mathfrak p\)-primary,
\(V_F(Q)=V_F(\mathfrak p)\).

\textbf{Proof.} If \(f^r\in I\), then at a zero of \(I\) one has
\(f(\eta)^r=0\); fields have no nonzero nilpotents, so \(f(\eta)=0\).
The reverse inclusion follows from \(I\subset\sqrt I\). Apply this to
\(\sqrt Q=\mathfrak p\). \(\square\)

If \(\mathfrak p_1,\ldots,\mathfrak p_s\) are the minimal primes over
\(I\), then

\[
V_\Omega(I)=\bigcup_i V_\Omega(\mathfrak p_i).
\]

Indeed \(\sqrt I=\bigcap_i\mathfrak p_i\). A point outside every
\(V(\mathfrak p_i)\) admits \(f_i\in\mathfrak p_i\) nonzero there, and
the product \(\prod f_i\in\sqrt I\) stays nonzero there, a
contradiction. By Theorem 2.1 these sets are the specializations of the
generic zeros, when those zeros are realized in \(\Omega\).

The strong Nullstellensatz implies

\[
I_k(V_\Omega(I))=\sqrt I.
\]

For a general \(k\) this follows by applying it to \(I\Omega[x]\) and
contracting: field extension is faithfully flat, and
\(I\Omega[x]\cap k[x]=I\); the same holds for radicals, since powers
detect them. This does not claim that \(\Omega\) need be large enough
for a generic zero: an algebraic closure of \(k\) already suffices for
the equality of ideals of zeros.

The primary ideal \((x^2,y)\) has only the origin as its geometric zero.
The quotient still contains a nonzero nilpotent class \(x\). A set of
field-valued zeros detects the radical, so multiplicity needs more
algebra.

\subsection{4. Counting parameters and excluding embedded
dimensions}\label{counting-parameters-and-excluding-embedded-dimensions}

Krull dimension is the supremum of the lengths of strict prime chains.
For a prime \(\mathfrak p\), the cited dimension theorem gives

\[
\dim(R/\mathfrak p)=\operatorname{trdeg}_k k(\xi).
\]

Noether's 1923 prime-chain formulation describes the same number as the
parameter formulation. The catenarity and height formula belong to the
cited commutative-algebra lesson. For the parabola, \((0)\subsetneq(t)\)
in \(k[t]\) is a length-one chain, and its function field has one
parameter.

An ideal in a polynomial ring is \textbf{unmixed} here if every
associated prime of its quotient has the same height. This includes
embedded associated primes in the test. Equidimensionality, which only
tests minimal primes, is weaker. For \((x^2,xy)\), the associated
heights are one and two, so it is not unmixed, although its reduced
closed set is just one line.

\textbf{Theorem 4.1 (Macaulay's unmixedness theorem).} In
\(k[x_1,\ldots,x_n]\), an ideal generated by \(r\) elements and having
height \(r\) is unmixed: every associated prime of its quotient has
height \(r\).

\textbf{Proof using the internal prerequisites.} Polynomial
localizations are regular by Proposition 3.3 of
\href{../prerequisites/AG-CA/AG-CA-14.html}{Regular local rings}. They
are Cohen--Macaulay by Theorem 6.1 of
\href{../prerequisites/AG-CA/AG-CA-12.html}{Regular sequences, depth and
Cohen--Macaulay modules}. Let \(\mathfrak q\) be associated to the
quotient, and localize at \(\mathfrak q\). Every minimal prime
\(\mathfrak p\) over \(I\) contained in \(\mathfrak q\) has height
exactly \(r\): height is at least \(r\) by the assumed height of \(I\),
and at most \(r\) by the generalized principal ideal theorem for \(r\)
generators. The polynomial height and catenarity formulas in the
dimension prerequisite give
\(\dim (R_{\mathfrak q}/IR_{\mathfrak q})=\operatorname{ht}\mathfrak q-r\).
Corollary 6.2 of the depth lesson makes this quotient Cohen--Macaulay
with that dimension. Since its maximal ideal is associated, it has depth
zero, by Lemma 2.2 of that same lesson. Its dimension is therefore zero,
forcing \(\operatorname{ht}\mathfrak q=r\). This proves unmixedness. The
generalized principal ideal theorem is proved in
\href{../prerequisites/AG-CA/AG-CA-11.html}{Dimension theory of
Noetherian local rings}, Section 3. \(\square\)

This argument uses regularity of polynomial localizations. The same
number-of-generators condition in an arbitrary Noetherian ring is
insufficient.

\subsection{5. When the variety survives every field
extension}\label{when-the-variety-survives-every-field-extension}

The prime \(\mathfrak p\) is \textbf{absolutely prime} if
\(\mathfrak p\bar k[x]\) is prime for an algebraic closure \(\bar k\) of
\(k\). A \(k\)-algebra is \textbf{geometrically integral} when its
tensor product with every extension field is a domain.

For a finitely generated field extension, \textbf{separable over \(k\)}
means \textbf{separably generated}: it has a transcendence basis
\(t_1,\ldots,t_d\) such that the finite algebraic extension over
\(k(t_1,\ldots,t_d)\) is separable. A \(k\)-algebra is
\textbf{geometrically reduced} when its tensor product with every
extension field is reduced. The next four lemmas provide all
field-extension facts used in the criterion, including the
imperfect-field case.

\textbf{Lemma 5.A (finite separable extensions).} A finite separable
extension \(E/F\) has a primitive element. For any extension field
\(H/F\), \(E\otimes_F H\) is a finite product of finite separable
extensions of \(H\), hence is reduced. In characteristic \(p>0\), a
finite extension \(E/F\) is separable if and only if \(E=FE^p\), where
\(FE^p\) denotes the subfield generated by \(F\) and all \(p\)-th powers
in \(E\).

\textbf{Proof.} First recall the root-count argument. In a finite tower
of simple algebraic extensions, each embedding of an intermediate field
into an algebraic closure extends by choosing a root of the transported
minimal polynomial. The number of choices is its degree when that
polynomial is separable and is at most its degree in general. Thus a
separable tower has exactly its total degree many embeddings. Every
element of its top field is separable: an inseparable element would have
fewer roots than its degree, and counting embeddings by restriction to
the field it generates would give fewer embeddings than the total
degree. This also proves transitivity of finite algebraic separability
and that the separable elements in a finite extension form a subfield.

If \(F\) is infinite and \(E=F(a,b)\) is separable, there are \([E:F]\)
embeddings. Two distinct embeddings disagree on \(a\) or \(b\), so their
images of \(a+cb\) coincide for at most one \(c\in F\). Avoiding the
finitely many forbidden values gives an element with \([E:F]\) distinct
conjugates. Its degree is therefore \([E:F]\), so it generates \(E\).
Induct on a finite list of generators. If \(F\) is finite, then \(E\) is
finite. Its multiplicative group is cyclic: for its exponent \(m\),
combining commuting elements of maximal prime-power orders produces an
element of order \(m\); all group elements are roots of \(T^m-1\), so
the group has at most \(m\) elements, and it has at least \(m\) because
of this element. A generator of that group generates the field. This
proves the primitive-element assertion in both cases.

Write \(E=F(a)\) with separable minimal polynomial \(f\). Bézout's
identity for \(f,f'\) remains valid over \(H\), so the irreducible
factors of \(f\) in \(H[T]\) are distinct and separable. The Chinese
remainder theorem gives

\[
E\otimes_F H=H[T]/(f)\cong\prod_j H[T]/(g_j),
\]

with each factor a finite separable field. The Chinese remainder theorem
here follows directly from the Bézout identities between distinct
irreducible polynomials: their ideals are comaximal, and the usual
remainder maps have kernel their product.

For the last assertion, if \(a\in E\) is separable over \(F\), its
minimal polynomial over \(F(a^p)\) is separable and divides \(T^p-a^p\),
which has only one distinct root. That minimal polynomial has degree
one, so \(a\in F(a^p)\). Consequently separability implies \(E=FE^p\).

Conversely let \(E_s\) be the subfield of elements of \(E\) separable
over \(F\). Every irreducible polynomial in characteristic \(p\) can be
written \(g(T^{p^e})\) with \(g'\ne0\); \(g\) is irreducible and hence
separable. It follows that some \(p\)-power of each element of \(E\)
belongs to \(E_s\). A finite generating list supplies one exponent \(N\)
such that \(E^{p^N}\subset E_s\). If \(E=FE^p\), then \(E=E_sE^p\), and
iteration gives \(E=E_sE^{p^N}=E_s\). This is separability. \(\square\)

\textbf{Lemma 5.B (separability and geometric reducedness).} For a
finitely generated field extension \(L/k\), the following are
equivalent:

\begin{enumerate}
\def\labelenumi{\arabic{enumi}.}
\tightlist
\item
  \(L/k\) is separably generated.
\item
  \(L\otimes_k K\) is reduced for every extension field \(K/k\).
\item
  \(L\otimes_k\bar k\) is reduced.
\end{enumerate}

In characteristic \(p>0\), they are also equivalent to reducedness of
\(L\otimes_k k^{1/p}\), where \(k^{1/p}=\{c\in\bar k:c^p\in k\}\).

\textbf{Proof.} Suppose \(F=k(t_1,\ldots,t_d)\subset L\) is a separating
transcendence basis. For any \(K/k\),

\[
B=F\otimes_k K=(k[t_1,\ldots,t_d]\setminus\{0\})^{-1}K[t_1,\ldots,t_d]
\]

is a domain. If \(H=\operatorname{Frac}(B)\), the injection
\(B\hookrightarrow H\) remains injective after tensoring with the
finite-dimensional \(F\)-vector space \(L\): choose an \(F\)-basis to
see this coordinate by coordinate. Hence

\[
L\otimes_k K=L\otimes_F B\hookrightarrow L\otimes_F H.
\]

The ring on the right is reduced by Lemma 5.A, so its subring on the
left is reduced. Thus (1) implies (2), which implies (3). In
characteristic zero, every finitely generated field extension has a
separating transcendence basis: choose any transcendence basis, and its
remaining finite extension is separable because an irreducible
nonconstant polynomial has nonzero derivative. This proves the
equivalences in that characteristic.

Assume now that the characteristic is \(p>0\) and that
\(L\otimes_k k^{1/p}\) is reduced. Condition (3) implies this
hypothesis, since tensoring the injection \(k^{1/p}\subset\bar k\) with
the \(k\)-vector space \(L\) is injective. Place the compositum
\(Lk^{1/p}\) in an algebraic closure of \(L\). The multiplication map

\[
L\otimes_k k^{1/p}\longrightarrow Lk^{1/p}
\]

is injective. Indeed, for \(z=\sum_i a_i\otimes c_i\),
\(z^p=(\sum_i a_i^pc_i^p)\otimes1\); if the multiplication image of
\(z\) vanishes, this power is zero, and reducedness gives \(z=0\).
Frobenius gives an isomorphism from this tensor ring to
\(L^p\otimes_{k^p}k\). Therefore its multiplication map into \(L\) is
also injective. Explicitly, any elements of \(k\) linearly independent
over \(k^p\) remain linearly independent over \(L^p\).

Take any transcendence basis \(t_1,\ldots,t_d\), put
\(F=k(t_1,\ldots,t_d)\), and write \(n=[L:F]\). Frobenius is an
isomorphism of the extensions \(L/F\) and \(L^p/F^p\), so
\([L^p:F^p]=n\). The injective tensor product \(k\otimes_{k^p}L^p\) is
free of rank \(n\) over its subring \(k\otimes_{k^p}F^p\). The fraction
field of the latter subring is \(kF^p=k(t_1^p,\ldots,t_d^p)\).
Localizing at its nonzero elements consequently gives a domain of
dimension \(n\) over \(kF^p\). A finite-dimensional domain over a field
is a field: multiplication by a nonzero element is an injective
endomorphism of a finite-dimensional vector space, hence surjective.
This field is the compositum \(kL^p\). Thus

\[
[kL^p:kF^p]=n,\qquad
[L:kL^p]=\frac{[L:F][F:kF^p]}{[kL^p:kF^p]}=p^d.
\tag{5.1}
\]

Here \([F:kF^p]=p^d\): the monomials \(t_1^{e_1}\cdots t_d^{e_d}\),
\(0\le e_i<p\), are independent by grouping polynomial exponents modulo
\(p\) after clearing denominators, and they span because a polynomial
denominator \(q(t)\) has inverse \(q(t)^{p-1}/q(t)^p\), with
\(q(t)^p\in kF^p\).

Put \(M=kL^p\). Each element of \(L\) has its \(p\)-th power in \(M\).
Whenever an element \(u\) is outside an intermediate field containing
\(M\), adjoining it has degree \(p\): its minimal polynomial divides
\((T-u)^p\), and a monic factor of degree \(e\) with \(0<e<p\) would
have coefficient \(-eu\), forcing \(u\) into that field. By (5.1),
adjoining such elements successively stops after exactly \(d\) steps. We
obtain

\[
L=M(u_1,\ldots,u_d),
\qquad
\{u_1^{e_1}\cdots u_d^{e_d}:0\le e_i<p\}
\text{ is an }M\text{-basis of }L.
\tag{5.2}
\]

These \(u_i\) are algebraically independent over \(k\). If not, choose a
nonzero relation \(f(u)=0\) of least total degree. Group its exponents
to write

\[
f(X)=\sum_{0\le e_i<p}X^e f_e(X_1^p,\ldots,X_d^p),\qquad f_e\in k[Y_1,\ldots,Y_d].
\]

The basis (5.2) forces every \(f_e(u_1^p,\ldots,u_d^p)\) to vanish.
Choose a nonzero \(f_e\), and express its finitely many coefficients
using elements \(b_1,\ldots,b_r\in k\) independent over \(k^p\):
\(f_e(Y)=\sum_j b_jg_j(Y)\), with \(g_j\in k^p[Y]\). The injectivity of
\(k\otimes_{k^p}L^p\to L\) forces each \(g_j(u_1^p,\ldots,u_d^p)=0\).
Taking \(p\)-th roots of its coefficients gives a polynomial
\(h_j\in k[X]\) with \(g_j(u_1^p,\ldots,u_d^p)=h_j(u)^p\). At least one
\(h_j\) is nonzero, and its total degree is at most
\(\deg f_e\le\deg f/p<\deg f\). Its vanishing contradicts minimality. If
\(d=0\), the asserted independence is empty and needs no argument.

The \(u_i\) therefore form a transcendence basis. The extension
\(L/k(u_1,\ldots,u_d)\) is finite, and (5.2) says
\(L=k(u_1,\ldots,u_d)L^p\). Lemma 5.A makes that finite extension
separable. This proves (1) from the \(k^{1/p}\) test, completing all
equivalences. \(\square\)

\textbf{Lemma 5.C (algebraic-closure tests).} For any \(k\)-algebra
\(S\),

\[
S\text{ is geometrically integral}
\Longleftrightarrow S\otimes_k\bar k\text{ is a domain},
\]

and

\[
S\text{ is geometrically reduced}
\Longleftrightarrow S\otimes_k\bar k\text{ is reduced}.
\]

No finite-generation hypothesis on \(S\) or on the test extension is
required.

\textbf{Proof.} Only the reverse implications need proof. Suppose
\(u,v\in S\otimes_k K\) are nonzero with \(uv=0\). Express them as
\(u=\sum_i a_i\otimes c_i\), \(v=\sum_j b_j\otimes d_j\), choosing the
\(a_i\) linearly independent over \(k\), and likewise the \(b_j\).
Choose \(c_{i_0},d_{j_0}\ne0\). Form the nonzero finite-type
\(k\)-algebra

\[
C=k[c_i,d_j,(c_{i_0}d_{j_0})^{-1}]\subset K.
\]

Expand the finitely many products \(a_ib_j\) in a \(k\)-basis of their
span. The equation \(uv=0\) says that the corresponding finite list of
polynomial expressions in the \(c_i,d_j\) vanishes in \(C\). Choose a
maximal ideal of \(C\). Its residue field is finite algebraic over
\(k\), by the proved weak Nullstellensatz, and embeds over \(k\) in
\(\bar k\). The resulting map \(C\to\bar k\) preserves all these
equations and keeps \(c_{i_0},d_{j_0}\) nonzero, since their product is
invertible in \(C\). Linear independence of the chosen lists then gives
nonzero images \(\bar u,\bar v\in S\otimes_k\bar k\) with
\(\bar u\bar v=0\), contradicting the domain hypothesis. The zero ring
cannot intervene: if \(S\otimes_k\bar k\) is a domain, then \(S\ne0\),
and tensoring its nonzero vector space with any field extension remains
nonzero.

For reducedness, take a nonzero \(u=\sum_i a_i\otimes c_i\) with
\(u^m=0\), choose \(c_{i_0}\ne0\), and use \(C=k[c_i,c_{i_0}^{-1}]\).
Expand the finitely many products arising in \(u^m\) in a basis of their
span. The same specialization produces a nonzero nilpotent in
\(S\otimes_k\bar k\), a contradiction. This proves the second test as
well. \(\square\)

\textbf{Lemma 5.D (purely inseparable scalar extension).} Let \(F/E\) be
a field extension and \(P/E\) a purely inseparable algebraic extension.
Then \(F\otimes_E P\) has exactly one prime ideal. It is a domain
whenever it is reduced.

\textbf{Proof.} In characteristic zero \(P=E\), so this is immediate. In
characteristic \(p>0\), choose compatible embeddings in an algebraic
closure of \(F\), and multiply tensors into the compositum \(FP\). For
any \(z=\sum_i a_i\otimes c_i\), there is \(N\) with every
\(c_i^{p^N}\in E\); hence \(z^{p^N}\) belongs to the embedded field
\(F\). If the multiplication image of \(z\) is zero, this power is zero;
conversely a nilpotent has zero image in a field. Thus the kernel is
exactly the nilradical. If the image of \(z\) is nonzero, its scalar
power is nonzero, and \(z^{p^N-1}/z^{p^N}\) is already an inverse in the
tensor ring. The image is therefore a field. Every prime contains the
nilradical, and its quotient is this field, proving uniqueness. If the
ring is reduced, that unique prime is zero, so the ring is a domain.
\(\square\)

\textbf{Theorem 5.1 (regular-extension criterion).} For
\(A=R/\mathfrak p\) and \(L=\operatorname{Frac}(A)\), the following are
equivalent:

\begin{enumerate}
\def\labelenumi{\arabic{enumi}.}
\tightlist
\item
  \(\mathfrak p\) is absolutely prime.
\item
  \(A\) is geometrically integral.
\item
  \(L/k\) is separable and every element of \(L\) algebraic over \(k\)
  belongs to \(k\).
\end{enumerate}

A field extension satisfying the third condition is called
\textbf{regular}.

\textbf{Proof.} The first two conditions are equivalent by Lemma 5.C,
since \(A\otimes_k\bar k=\bar k[x_1,\ldots,x_n]/\mathfrak p\bar k[x]\).
For every extension \(K/k\), tensoring the injection
\(A\hookrightarrow L\) over the field \(k\) gives an injection
\(A\otimes_k K\hookrightarrow L\otimes_k K\). Also \(L\otimes_k K\) is
the localization of \(A\otimes_k K\) at the images of
\(A\setminus\{0\}\). These images are nonzero because field extension
preserves the injection of \(A\). Hence these two rings are domains
simultaneously: one implication uses the injection, the other uses
localization of a domain at nonzero elements. It suffices to establish
the criterion for \(L\).

If \(L\otimes_k\bar k\) is a domain, it is reduced, so Lemma 5.B gives
separability of \(L/k\). If \(\alpha\in L\) is algebraic over \(k\),
tensoring the subfield injection gives
\(k[\alpha]\otimes_k\bar k\hookrightarrow L\otimes_k\bar k\). The source
is \(\bar k[T]/(f)\), where \(f\) is the minimal polynomial of
\(\alpha\). If \(\deg f>1\), write \(f=(T-a)g\) in \(\bar k[T]\). Both
factors have positive degree smaller than \(\deg f\), so both give
nonzero classes whose product is zero. This is impossible in a subring
of a domain, including when all roots of \(f\) coincide. Hence
\(\deg f=1\), giving \(\alpha\in k\).

Conversely assume the third condition. Any irreducible polynomial
\(f\in k[T]\) remains irreducible over \(L\). Indeed, the coefficients
of a monic factor in \(L[T]\) are elementary symmetric expressions in a
selection of roots in an algebraic closure of \(L\), so are algebraic
over \(k\). The constants hypothesis puts those coefficients in \(k\);
monic division then puts the complementary factor in \(k[T]\) too,
contradicting irreducibility. For a finite separable extension \(E/k\),
Lemma 5.A supplies a primitive element, and this irreducibility makes
\(L\otimes_k E\) a field. The separable closure \(k^{\mathrm{sep}}\) is
the directed union of its finite separable subextensions. Their tensor
rings inject into one another and are fields, so their union
\(F=L\otimes_k k^{\mathrm{sep}}\) is a field.

The extension \(\bar k/k^{\mathrm{sep}}\) is purely inseparable: in
characteristic \(p\), the minimal polynomial of any algebraic element
has the form \(g(T^{p^e})\) with \(g\) separable, so its \(p^e\)-th
power belongs to \(k^{\mathrm{sep}}\). In characteristic zero this stage
is absent. Consequently

\[
L\otimes_k\bar k=F\otimes_{k^{\mathrm{sep}}}\bar k
\]

has one prime by Lemma 5.D. Separability of \(L/k\) makes it reduced by
Lemma 5.B. Lemma 5.D therefore makes it a domain, proving absolute
primality and the equivalence of all three conditions. \(\square\)

\subsection{6. Two distinct failures}\label{two-distinct-failures}

Over \(\mathbb R\), the polynomial \(x^2+y^2\) is irreducible: a
degree-two factorization would be into homogeneous linear factors, but a
real line cannot lie in its zero set, which is only the origin. Since
\(\mathbb R[x,y]\) is a unique factorization domain, the generated ideal
is prime. Over \(\mathbb C\),

\[
x^2+y^2=(x+iy)(x-iy),
\]

so it is not absolutely prime. In its function field, \(y\ne0\) and
\((x/y)^2=-1\), yielding a new algebraic constant. Separability holds in
characteristic zero; algebraic closedness of the constants fails.

Over \(k=\mathbb F_p(t)\), the polynomial \(X^p-t\) is irreducible. The
element \(t\) is not a \(p\)-th power, as its order of vanishing at
\(t=0\) is one. If \(a^p=t\) in an algebraic closure, a monic proper
factor of \((X-a)^p\) would be \((X-a)^e\), \(0<e<p\); its coefficient
\(-ea\) would force \(a\in k\). This proves irreducibility. Its quotient
is the field \(k(t^{1/p})\). Over \(\bar k\) it becomes
\((X-t^{1/p})^p\); the class of \(X-t^{1/p}\) is nonzero by monic
division and is nilpotent. The field extension is inseparable and has
new algebraic constants, so both regularity conditions fail.

Separability cannot be dropped even if the constants condition holds.
Here is a complete example. Let \(s,x,y\) be independent indeterminates
over \(\mathbb F_p\), put

\[
L=\mathbb F_p(s,x,y),\qquad t=y^p-sx^p,\qquad k=\mathbb F_p(s,t).
\tag{6.1}
\]

The elements \(s,x,t\) are algebraically independent: in a polynomial in
\(t\), substitution of \(y^p-sx^p\) makes its highest \(t\)-degree term
the unique highest \(y\)-degree term. Thus \(x\) is transcendental over
\(k\). The element \(sx^p+t\) is not a \(p\)-th power in \(k(x)\), since
its formal derivative with respect to \(t\) is one, whereas every
\(p\)-th power has derivative zero. The preceding monic-factor argument
proves that

\[
A=k[x,y]/(y^p-sx^p-t)
\]

is a domain with fraction field \(L\). Over \(\bar k\), its equation is
\((y-s^{1/p}x-t^{1/p})^p\). Write \(h=y-s^{1/p}x-t^{1/p}\). The ring
\(A\otimes_k\bar k=\bar k[x,y]/(h^p)\) has nilradical \((h)\), quotient
\(\bar k[x]\), and a nonzero nilpotent \(h\). Reduction modulo \(h\)
embeds \(A\) in \(\bar k[x]\): the displayed value of \(y\) has the
irreducible degree-\(p\) polynomial just proved as its minimal
polynomial over \(k(x)\). Thus the elements of \(A\setminus\{0\}\) stay
outside \((h)\). Localize at those elements to obtain
\(L\otimes_k\bar k\). Its reduction is a subring of \(\bar k(x)\)
containing \(k(x)\), because every nonzero polynomial in \(k[x]\) is
among the denominators. Every element of this subring is algebraic over
\(k(x)\): only finitely many coefficients in the algebraic extension
\(\bar k/k\) occur in each rational function. A nonzero such element has
a minimal polynomial with nonzero constant coefficient, which expresses
its inverse as a polynomial in it with coefficients in \(k(x)\). Hence
the reduction is a field. Since \(h\) is nilpotent, this makes \((h)\)
the unique prime after localization. Also \(h\ne0\) there: its
annihilator before localization is \((h^{p-1})\subset(h)\), disjoint
from the denominators. Thus \(L\otimes_k\bar k\) is nonreduced, so
\(L/k\) is not separably generated by Lemma 5.B.

Nevertheless \(k\) is algebraically closed in \(L\). First there are no
new separable algebraic constants. A proper finite separable subfield
\(E/k\) would give an injection of \(E\otimes_k\bar k\), a product of at
least two copies of \(\bar k\), into \(L\otimes_k\bar k\). A ring with
one prime has no idempotents except zero and one: its unique prime is
its nilradical, so for an idempotent \(e\), either \(e\) or \(1-e\) is
both nilpotent and idempotent and must be zero. This rules out that
injection.

To exclude purely inseparable constants, put \(U=x^p,V=y^p\). Then
\(L^p=\mathbb F_p(s^p,U,V)\). If \(f\in k\cap L^p\), differentiate in
\(\mathbb F_p(s,U,V)\), holding \(U,V\) fixed. Since \(t=V-sU\), the
result is

\[
0=\frac{\partial f}{\partial s}-U\frac{\partial f}{\partial t},
\]

where the two derivatives on the right are computed in
\(k=\mathbb F_p(s,t)\). The element \(U\) is transcendental over \(k\),
so both derivatives vanish. The monomials \(s^it^j\), \(0\le i,j<p\),
are a basis of \(k\) over \(k^p\), by the same rational monomial
argument used in (5.1). Differentiating their expansion first in \(s\)
and then in \(t\) forces all coefficients except the constant one to
vanish. Hence \(f\in k^p\), proving \(k\cap L^p=k^p\). If \(a\in L\) has
\(a^p\in k\), then \(a^p=c^p\) for some \(c\in k\), and uniqueness of
\(p\)-th roots in a field gives \(a=c\). Repetition excludes every
nontrivial purely inseparable constant. Finally, an arbitrary element
algebraic over \(k\) has a suitable \(p\)-power separable over \(k\), by
the polynomial argument in Lemma 5.A. That power is in \(k\), and the
purely inseparable argument puts the element itself in \(k\). Thus (6.1)
satisfies the constants condition but fails separability, as claimed.

\subsection{7. Exercises}\label{exercises}

\begin{enumerate}
\def\labelenumi{\arabic{enumi}.}
\tightlist
\item
  \textbf{Basic.} Find a generic zero for \((y^2-x^3)\), and describe
  its specializations over an algebraically closed extension.
\item
  \textbf{Intermediate.} Show that a point whose coordinates are
  algebraic over \(k\) specializes the generic zero of \(\mathfrak p\)
  exactly when it lies in \(V(\mathfrak p)\). Explain why a map on
  coordinate rings suffices.
\item
  \textbf{Intermediate.} Prove that \((x^2+y^2)\) is prime over
  \(\mathbb R\) and compute its complex extension as an intersection of
  prime ideals.
\item
  \textbf{Intermediate.} For \((X^p-t,Y)\subset\mathbb F_p(t)[X,Y]\),
  identify the quotient field and its base change to an algebraic
  closure.
\item
  \textbf{Advanced.} Prove the regular-extension criterion in
  characteristic zero, identifying exactly where the characteristic-zero
  hypothesis removes a step.
\end{enumerate}

\subsection{8. Solutions}\label{solutions}

\textbf{1.} Use \((t^2,t^3)\), with \(t\) transcendental. Dividing by
the monic relation in \(y\) leaves \(a(x)+yb(x)\). Its evaluation has
disjoint even powers from \(a(t^2)\) and odd powers at least three from
\(t^3b(t^2)\), so vanishing forces both polynomials zero. The kernel is
exactly the given ideal, which is therefore prime. Every cusp point is
parametrized: if \(x\ne0\), set \(t=y/x\), giving \(t^2=x,t^3=y\); if
\(x=0\), then \(y=0\) and set \(t=0\). These are all specializations by
Theorem 2.1.

\textbf{2.} Containment of relation ideals is exactly the condition that
every polynomial in \(\mathfrak p\) vanish at the point; this is
membership in \(V(\mathfrak p)\). The quotient map therefore descends to
the coordinate algebra. Algebraicity ensures the point's coordinate
algebra is a finite field extension, by adjoining finitely many
algebraic elements, but it does not make a specialization map extend to
the generic function field.

\textbf{3.} The real irreducibility argument in Section 6 gives
primality. In \(\mathbb C[x,y]\) the distinct irreducible linear factors
are relatively prime as elements, so their principal ideals intersect in
their product:

\[
(x^2+y^2)=(x+iy)\cap(x-iy).
\]

These two ideals are not comaximal: their sum is \((x,y)\). The complex
variety is two lines meeting at the origin.

\textbf{4.} The quotient is \(\mathbb F_p(t)(t^{1/p})\), a purely
inseparable degree-\(p\) extension. The base-changed quotient is
\(\bar k[\epsilon]/(\epsilon^p)\) with \(\epsilon=X-t^{1/p}\) and
\(Y=0\). It is nonreduced and not a domain. The extension is not
separable and \(k\) is not algebraically closed in it.

\textbf{5.} In characteristic zero every finitely generated field
extension is separably generated. If \(k\) is algebraically closed in
\(L\), the monic-factor argument in Theorem 5.1 keeps every finite
separable minimal polynomial irreducible over \(L\), so
\(L\otimes_k\bar k\) is the directed union of fields and is a domain.
There is no purely inseparable stage. Conversely any algebraic
\(\alpha\in L\setminus k\) would give a degree-greater-than-one
separable polynomial whose tensor quotient over \(\bar k\) is a product
of fields, injecting into \(L\otimes_k\bar k\); this contradicts its
being a domain. The passage between \(A\) and \(L\) is the injective
localization argument already proved.

\subsection{In Noether's words}\label{in-noethers-words}

\href{https://zenodo.org/records/21923146/files/00_EMMY_NOETHER_EN_COMPLETE_LINKED_READER.pdf}{Elimination
Theory and General Ideal Theory, work 24} announces the arithmetic
dimension in its introduction and develops zeros in Sections 2--3,
dimension in Section 4, and absolute prime ideals in Section 7. Thus
absolute primality occurs later than the initial zero theory. The short
announcements are works 18 and 25.
\href{https://zenodo.org/records/21908301/files/10-de.pdf}{The German
authority} gives the original text.

\subsection{References}\label{references}

\begin{itemize}
\tightlist
\item
  Emmy Noether,
  \href{https://zenodo.org/records/21908301/files/10-de.pdf}{\emph{Eliminationstheorie
  und allgemeine Idealtheorie}, work 24 in the freely readable German
  edition}, Mathematische Annalen \textbf{90} (1923), 229--261,
  introduction and Sections 1--4, 7;
  \href{https://zenodo.org/records/21923146/files/00_EMMY_NOETHER_EN_COMPLETE_LINKED_READER.pdf}{freely
  readable English edition}, work 24.
\item
  The Stacks project authors, freely accessible comparison statements:
  \href{https://stacks.math.columbia.edu/tag/030I}{separable extensions,
  Tag 030I}; \href{https://stacks.math.columbia.edu/tag/0322}{geometric
  reducedness, Tag 0322};
  \href{https://stacks.math.columbia.edu/tag/037P}{algebraically closed
  constants, Tag 037P};
  \href{https://stacks.math.columbia.edu/tag/0FWF}{algebraic-closure
  test for geometric integrality, Tag 0FWF}. The source retains its GFDL
  licence. Lemmas 5.A--5.D and Theorem 5.1 above supply the complete
  proofs needed here, including the finitely generated separability
  convention and arbitrary test extensions.
\end{itemize}

\subsection{Editable sources}\label{editable-sources}

\href{ideal-theory-in-rings-source.zip}{Complete source archive} ·
\href{ideal-theory-in-rings.tex}{Complete course in LaTeX} ·
\href{NOE-IDL-04.tex}{This lesson in LaTeX}. The archive includes all
five lesson texts, the figures and their reproducible source.

\end{document}
